\chapter{Generative AI Prompts}
\label{app:genai}

This appendix contains prompts used with Claude Code to generate diagrams and code samples in this thesis. These prompts can be used to regenerate or modify the artifacts.

\section{Architecture Diagram Prompt}

The following prompt was used to generate the TikZ architecture diagram in Figure~\ref{fig:architecture}:

\begin{lstlisting}[frame=single, backgroundcolor=\color{backcolour}, rulecolor=\color{framecolour}, basicstyle=\small\ttfamily, breaklines=true]
Create a TikZ architecture diagram for a data lakehouse security
framework with these specifications:

Layout: Vertical flow, top-to-bottom

Components (in order):
1. Top row: 5 external source boxes horizontally: GitHub,
   SonarQube, Snyk, NVD, EPSS
2. Ingestion: "dlt Connectors" box (orange tint)
3. Bronze layer: Blue background strip containing 3 yellow boxes
   (github.*, sonarqube.*, nvd.*/epss.*)
4. Silver layer: Blue background strip containing 3 gray boxes
   (dimensions, facts, relationships)
5. Gold layer: Blue background strip containing 2 green boxes
   (ml_enriched, metrics)
6. Wrap Bronze/Silver/Gold in a larger box labeled "Databricks"
7. Below Databricks: Lakebase (ops.views), FastAPI Backend,
   React Frontend

Arrows: Connect layers with labeled arrows showing data flow:
- Sources to dlt Connectors
- dlt Connectors to Bronze
- Bronze to Silver (label: DLT)
- Silver to Gold (label: DLT + ML)
- Gold to Lakebase (label: Materialization)
- Lakebase to Backend (label: SQL)
- Backend to Frontend (label: REST API)

Style: Use soft colors, rounded corners, small fonts for
compact display.
\end{lstlisting}

